\section{Conclusion}
\label{sec:conclusion}

In a two-country world, the bilateral exchange rate, the effective exchange rate and the terms of trade coincide, and a tariff appreciates the tariffing country's currency wherever Marshall--Lerner holds. In a multilateral world they separate. We have characterized what the real trade channel implies once they do.

The general response is the inverse of a multilateral adjustment matrix applied to a vector of tariff disturbances. A multilateral Marshall--Lerner condition makes adjustment well behaved without fixing the sign of any bilateral response. Under symmetry, average protection moves the effective rate and relative treatment moves bilateral rates around it. In a nested-CES model the conventional result survives exactly at the effective level under unilateral tariffs. It fails bilaterally against bystanders when cross-origin substitution exceeds $N$ times the domestic diversion margin, and against everyone under retaliation once it exceeds the domestic diversion margin itself. An exact three-country representation shows that the role of $N$ comes from the inward-looking, highly substitutable trade of the bystander bloc, not from its size. Under asymmetry, even the effective and welfare-relevant conclusions can fail: tariffing a small, lightly weighted partner whose goods are easily replaced can worsen the tariffing country's aggregate terms of trade.

The implications run both ways. Appreciation against a tariffed partner is not evidence of broad currency strength. A multilateral depreciation during a discriminatory or retaliatory tariff episode cannot be attributed to financial channels on the basis of the two-country prediction alone. Conversely, when tariffs are nearly uniform, the trade channel predicts effective appreciation, and observed weakness must come from elsewhere.

Natural next steps are to embed the decomposition in a quantitative multi-country model with input--output linkages and heterogeneous elasticities, and to add nominal rigidities, invoicing and a financial account. This would let the long-run real benchmark derived here interact with the short-run monetary and portfolio dynamics that dominate exchange rates at business-cycle frequencies.
